\section{Discussion}
\label{sec:scope}

The query improvement comes from combining two ideas.  First, the algorithm
stores each candidate set as a conjunction of the adjacency constraints
created by earlier pivots.  This representation supports quantum sampling
without enumerating the surviving vertices.  Second, the scale-aware analysis
allocates estimation error according to the cost of the recursion level.  The
resulting schedule keeps the cumulative loss in candidate-set size bounded
while concentrating most of the sampling effort where it is cheapest.

A natural next step is to determine the randomized classical query
complexity in the range between \cref{prop:classical-lower} and
\cref{thm:main}.  The implicit-set method may also extend to off-diagonal or
hypergraph Ramsey search.  Quantum walks already give sublinear-query
algorithms for fixed sub-hypergraph containment
\cite{LeGallNishimuraTani2016}; a Ramsey algorithm would additionally need to
exploit the totality and recursive structure of the problem.  For multiple
edge colours, \cref{cor:multicolour} gives a size-biased algorithm, while a
scale-aware analogue could yield a sharper dependence on the recursion depth.

Appendix~\ref{app:experiments} reports finite experiments that illustrate the
survival recurrence and the conditional uniformity of the sampler.  The main
query bound follows from the analytic argument in
\cref{sec:model,sec:algorithm,sec:correctness,sec:complexity}.

The constructive diagonal Ramsey recursion admits a quantum implementation
over implicit candidate sets.  Capped uniform search and a scale-aware
concentration schedule give $O(2^K K\log(K/\eta))$ edge queries, or
$\widetilde O(\sqrt N)$ queries at the succinct Ramsey scale.
